\section{Problema e bibliotecas}

\begin{frame}[t,label=cd-problema]{O que é Ciência de Dados?}
Ciência de Dados combina programação, Matemática e Estatística para analisar dados e apoiar decisões.
\begin{EMASCnote}
\textbf{Problema de estudo:} uma equipe acadêmica acompanha as notas de três cursos. Em quais há notas pendentes? Como se distribuem as notas disponíveis?
\end{EMASCnote}
\smallskip
\textbf{Base de dados:} 18 registros fictícios. Alguns exemplos:
\begin{tabularx}{\linewidth}{@{}lXrrr@{}}
\toprule
\texttt{id} & \texttt{curso} & \texttt{horas\_estudo} & \texttt{frequencia} & \texttt{nota} \\
\midrule
1 & Matematica & 2,0 & 0,70 & 5,9 \\
6 & Matematica & 3,5 & 0,80 & \\
\bottomrule
\end{tabularx}
\smallskip

\small
O ID identifica o estudante. A frequência é a proporção de aulas frequentadas (0,70 = 70\%). O campo vazio indica nota ainda não registrada neste exemplo.
\end{frame}

\begin{frame}[t]{Etapas de uma análise}
\begin{enumerate}
  \item \textbf{Entender o problema:} definir a pergunta que os dados deverão ajudar a responder.
  \item \textbf{Obter os dados:} reunir as informações necessárias à análise.
  \item \textbf{Preparar os dados:} verificar tipos, valores ausentes e inconsistências.
  \item \textbf{Explorar os dados:} calcular medidas estatísticas e construir gráficos.
  \item \textbf{Interpretar os resultados:} identificar padrões e reconhecer limitações.
  \item \textbf{Propor uma ação:} utilizar os resultados para justificar uma decisão e acompanhar sua aplicação.
\end{enumerate}
\end{frame}

\begin{frame}[fragile,t]{Bibliotecas de Python}
Uma biblioteca reúne funções e estruturas prontas para uso.

\begin{tabularx}{\linewidth}{@{}lX@{}}
\toprule
\textbf{Biblioteca} & \textbf{Aplicação} \\
\midrule
NumPy & Operações com vetores, matrizes e cálculos numéricos. \\
pandas & Leitura, manipulação e análise de dados tabulares. \\
Matplotlib & Construção de gráficos e visualização de dados. \\
\bottomrule
\end{tabularx}
\EMASCcodelabel{Importação}
\begin{lstlisting}
import numpy as np
import pandas as pd
import matplotlib.pyplot as plt
\end{lstlisting}
Os nomes \texttt{np}, \texttt{pd} e \texttt{plt} são abreviações das bibliotecas.
\end{frame}
